% Fichier généré automatiquement par tools/generate_corrections.py.
% Source : TEC/TEC-05/49_ElevateurBateaux/49_ElevateurBateaux.tex
% Statut : complet (2/2 question(s) renseignée(s), 0 reprise(s)).
\begin{corrigebox}[Corrigé extrait de la banque]
\CorrigeQuestion{1}
\begin{marginfigure}
\includegraphics[width=.5\linewidth]{49_02}
\end{marginfigure}
\CorrigeQuestion{2}
On isole l'ensemble $E=\{S2 ; T2,\}$. On applique le théorème de l’énergie cinétique à l’ensemble en mouvement dans le référentiel terrestre galiléen : 
	 $\mathcal{P}_{\text{int}}(E)+\pext{\overline{E}}{E}{R_g} = \dfrac{\dd \ec{E}{R_g} }{\dd t}$.

\textbf{Calcul des puissances externes}

$\pext{\text{pes}}{2}{R_g} = $

\textbf{	Calcul des puissances internes}
 $\mathcal{P}_{\text{int}}(E)=0$ car pas de frottement dans la liaison pivot.
\end{corrigebox}
